\section{Threats to Validity and Future Work}
\label{sec:limitations}

The thesis is intentionally bounded. It analyzes a validated set of
supervisor-provided artifacts rather than running a larger benchmark. This
improves auditability, but it limits the strength and generality of the
claims.

\subsection{Internal Validity}

The main internal validity constraint is the single-output design. Every
validated per-cell JSON sidecar used in the thesis reports \texttt{n\_runs =
1}. The results therefore do not estimate stochastic output variance,
confidence intervals, or repeated-run stability. The supervisor notebook
contains repeated-run and self-consistency machinery, but those sections are
not supported by the primary result files used here.

Prompt sensitivity is a second internal threat. Some prompts produce useful
graphs, some overpredict false positives, and some produce empty outputs for
DeepSeek. The reported ranking is therefore the observed result of specific
prompt templates, model versions, decoding settings, and parser rules, not a
universal ordering of prompt families.

\subsection{Construct Validity}

The ground-truth DAG is a physics-informed evaluation target. It is not a
formally identified causal model from interventions, randomized experiments,
or full system identification. The thesis therefore evaluates agreement with a
defined reference graph, not discovery of true causality in the world.

The parser also shapes the construct being measured. It enforces a fixed
feature vocabulary and arrow-based relationship syntax. A semantically useful
relationship expressed with an unsupported variable name or prose description
may be discarded, while a syntactically valid but conceptually weak edge is
counted as a prediction. Future work should test parser sensitivity,
schema-constrained generation, and JSON-only output validation.

\subsection{External Validity}

All primary results use one Tesla EV dataset and one six-feature graph. The
feature set is meaningful for energy-consumption and motion analysis, but it
does not cover all telematics variables used in fleet monitoring. Results may
differ for combustion vehicles, hybrid vehicles, richer EV signals,
maintenance records, GPS traces, or multi-vehicle fleet histories.

The model set is also limited to two local 8B-class open-source models. The
findings are relevant to local deployment scenarios, but they should not be
generalized to larger closed models, fine-tuned models, or tool-augmented
causal-discovery systems without additional experiments.

\subsection{Reproducibility Constraints}

The evidence is reproducible at the artifact level: the reported values can
be traced to saved CSV, JSON, and matrix files. It is not a full rerun-based
reproducibility study. A stronger benchmark would rerun each model-strategy
cell multiple times under fixed decoding settings, preserve every raw output,
and report edge probabilities alongside single-run graph metrics.

Because repeated outputs are absent from the validated artifact set, the
thesis cannot claim that self-consistency improves graph recovery. A proper
self-consistency experiment would need multiple independent outputs for the
same model, strategy, and dataset fingerprint, followed by a documented voting
or aggregation rule.

\subsection{Operational Governance}

The thesis does not use governance risk as a primary metric, but telemetry
governance remains relevant. A graph-recovery system may expose operational
patterns if prompts include unnecessary raw data or if outputs are published
without review. Future deployments should combine graph evaluation with data
scoping, output governance, and access controls.

\subsection{Future Work}

The next experimental step is a repeated-run benchmark over the same Tesla EV
task. Each model-strategy cell should be sampled multiple times, with edge
probabilities and self-consistency voting reported separately from single-run
metrics. A second step is cross-dataset validation: apply the same pipeline to
additional vehicle datasets and compare whether the same prompt strategies
remain useful.

Future work should also compare LLM graph proposals with conventional
statistical methods and expert-authored graphs. Such a comparison would
clarify whether LLMs add useful domain priors, merely echo feature-name
semantics, or overgenerate plausible but unsupported relationships.

\summarynote{Discussion takeaway}{The current evidence is strongest as a
controlled artifact-based comparison. It supports prompt decomposition as a
useful but unreliable graph-recovery aid, not as a validated causal-discovery
method.}
